\chapter{Background}

\comment{real world and eventually related work}

\section{Ecoinformatics}

\section{Biological study context}

An internal DISTAV methods document describes the biological study supported by the computational workflow. The study examines how climate change may alter the suitability of Alpine grassland habitats. It considers ten dry, alpine and subalpine habitat types from the EUNIS classification and 167 diagnostic plant species. Occurrence records are associated with environmental predictors at a spatial resolution of approximately 1~km. The predictors describe climate, soil and topographic heterogeneity. Species Distribution Models combine five modelling techniques and ensemble forecasts to estimate present and future habitat suitability.

The biological protocol also defines the analyses that follow model projection: changes in suitability, spatial shifts, habitat-level hotspot and cold-spot maps, and the relationship between suitability changes and species tolerance. This document provides the scientific context for the computing work, but it is not treated as the exact specification of every run. Dataset sizes, model repetitions and available climate scenarios changed during the project; the computational chapters therefore report the configuration recorded by the corresponding scripts and run manifests.

\section{Data used}

\section{Related work}
